% ============================================================
\chapter{Comparison Theorems}
\label{ch:comparison}
% ============================================================

\section{Introduction and motivation}

One of the most fruitful ideas in Riemannian geometry rests on a remarkably simple principle: if one knows that the curvature of a manifold is bounded, then its geometry cannot deviate too far from that of a ``model'' space of constant curvature. This comparison programme, initiated by Harry Rauch in the 1950s and then developed by Marcel Berger, Wilhelm Klingenberg, Jeff Cheeger, and Detlef Gromoll over the following decades, has produced some of the deepest results in global geometry. Rauch's theorem compares the behaviour of Jacobi fields; Bishop--Gromov compares volumes; Toponogov compares triangles. Each extracts global geometric conclusions from a local hypothesis on curvature, illustrating the power of curvature as a controlling invariant.

\section{Jacobi equation and model solutions}

Recall that along a unit-speed geodesic $\gamma$, an orthogonal Jacobi field
$J \perp \dot\gamma$ satisfies:
\[
    J'' + K(t)\, J = 0,
\]
where $K(t)$ denotes the sectional curvature of the plane $(\dot\gamma, J)$.

\begin{definition}[Model solutions]
For $\kappa \in \R$, the solution of $f'' + \kappa f = 0$ with $f(0) = 0$,
$f'(0) = 1$ is:
\[
    \operatorname{sn}_\kappa(t) = \begin{cases}
    \frac{1}{\sqrt{\kappa}}\sin(\sqrt{\kappa}\, t) & \text{if } \kappa > 0, \\
    t & \text{if } \kappa = 0, \\
    \frac{1}{\sqrt{-\kappa}}\sinh(\sqrt{-\kappa}\, t) & \text{if } \kappa < 0.
    \end{cases}
\]
\end{definition}

\begin{keyformulas}[title=\textbf{Comparison functions}]
\begin{align*}
    \operatorname{sn}_\kappa(t) &= \text{solution of } f'' + \kappa f = 0, \; f(0)=0, \; f'(0)=1, \\
    \operatorname{cn}_\kappa(t) &= \operatorname{sn}_\kappa'(t) = \text{solution of }
    f'' + \kappa f = 0, \; f(0)=1, \; f'(0)=0.
\end{align*}
\end{keyformulas}

\section{Sturm comparison lemma}

\begin{lemma}[Sturm]
\label{lem:sturm}
Let $f$ and $g$ be solutions of:
\[
    f'' + K_1(t)\, f = 0, \qquad g'' + K_2(t)\, g = 0,
\]
with $f(0) = g(0) = 0$, $f'(0) = g'(0) = 1$, and $K_1(t) \geq K_2(t)$.

Then $f(t) \leq g(t)$ on the interval $[0, t_1]$ where $t_1$ is the first zero
of $f$.
\end{lemma}

\begin{proof}
Consider $h = f'g - fg'$. We have $h(0) = 0$ and:
\[
    h' = f''g - fg'' = -K_1 fg + K_2 fg = (K_2 - K_1)fg \leq 0,
\]
as long as $f \geq 0$ and $g \geq 0$. Thus $h \leq 0$, i.e.\ $f'/f \leq g'/g$
(when $f, g > 0$). Integrating: $\ln f - \ln g$ is decreasing, and since
$\lim_{t \to 0^+} f(t)/g(t) = 1$, we get $f \leq g$.
\end{proof}

\section{Rauch comparison theorem}

\begin{theorem}[Rauch]
\label{thm:rauch}
Let $(M, g)$ and $(\bar{M}, \bar{g})$ be Riemannian manifolds, and
$\gamma : [0, L] \to M$, $\bar\gamma : [0, L] \to \bar{M}$ unit-speed geodesics
without conjugate points. Suppose:
\[
    K_M(\dot\gamma, \cdot) \leq K_{\bar{M}}(\dot{\bar\gamma}, \cdot).
\]
If $J$ and $\bar{J}$ are orthogonal Jacobi fields along $\gamma$ and $\bar\gamma$
with $J(0) = \bar{J}(0) = 0$ and $\norm{J'(0)} = \norm{\bar{J}'(0)}$, then:
\[
    \norm{J(t)} \geq \norm{\bar{J}(t)}, \quad \forall\, t \in [0, L].
\]
\end{theorem}

\begin{proof}[Proof sketch]
Consider the ratio $u(t) = \norm{J(t)}/\operatorname{sn}_\kappa(t)$ where
$\kappa$ is the lower curvature bound. By the Sturm lemma, compare $\norm{J}$
with the model solution. The key is the inequality:
\[
    \frac{d}{dt}\frac{\norm{J}'}{\norm{J}} \leq -K(t),
\]
which follows from the Jacobi equation and the Cauchy--Schwarz inequality.
The conclusion follows from Sturm comparison.
\end{proof}

\begin{corollary}[Comparison with model space]
If $K_M \leq \kappa$ (upper bound):
\[
    \norm{J(t)} \geq \operatorname{sn}_\kappa(t)\, \norm{J'(0)}.
\]
If $K_M \geq \kappa$ (lower bound):
\[
    \norm{J(t)} \leq \operatorname{sn}_\kappa(t)\, \norm{J'(0)}.
\]
\end{corollary}

\section{Toponogov's theorem}

\begin{theorem}[Toponogov -- global version]
\label{thm:toponogov}
Let $(M, g)$ be a complete Riemannian manifold with $K \geq \kappa$. Let
$\triangle(p, q, r)$ be a geodesic triangle in $M$ (formed by three minimizing
geodesics). Let $\bar\triangle(\bar{p}, \bar{q}, \bar{r})$ be a comparison
triangle in $M_\kappa$ (same side lengths). Then:

\textbf{Angle version:} The angles of the triangle in $M$ are \emph{greater
than or equal to} the corresponding angles of the comparison triangle:
\[
    \angle_p(q, r) \geq \bar\angle_{\bar{p}}(\bar{q}, \bar{r}).
\]

\textbf{Distance version:} For any point $x$ on the side $[q, r]$ and the
corresponding point $\bar{x}$ on $[\bar{q}, \bar{r}]$ (same proportions):
\[
    d(p, x) \leq d(\bar{p}, \bar{x}).
\]
\end{theorem}

\begin{warning}[title=\textbf{Hypotheses for Toponogov}]
Toponogov's theorem requires:
\begin{enumerate}
    \item Completeness of $(M, g)$.
    \item A \emph{lower} bound on sectional curvature.
    \item The sides of the triangle are \emph{minimizing} geodesics.
    \item If $\kappa > 0$, the perimeter does not exceed $2\pi/\sqrt{\kappa}$.
\end{enumerate}
\end{warning}

\begin{corollary}
If $K \geq 0$ (nonnegative curvature), geodesic triangles are ``fatter'' than
the corresponding Euclidean triangles: angles are larger and ``medians'' are shorter.
\end{corollary}

\section{Bishop--Gromov comparison theorem}

\begin{theorem}[Bishop--Gromov]
\label{thm:bishop_gromov}
Let $(M^n, g)$ be a complete Riemannian manifold with
$\operatorname{Ric} \geq (n-1)\kappa$. Let $V_\kappa(r)$ be the volume of the
ball of radius $r$ in the model space $M^n_\kappa$. Then the function:
\[
    r \mapsto \frac{\operatorname{Vol}(B(p, r))}{V_\kappa(r)}
\]
is nonincreasing for $r > 0$, and tends to $1$ as $r \to 0$.

In particular:
\[
    \operatorname{Vol}(B(p, r)) \leq V_\kappa(r), \quad \forall\, r > 0.
\]
\end{theorem}

\begin{proof}[Proof sketch]
In geodesic polar coordinates:
$\operatorname{Vol}(B(p,r)) = \int_{S^{n-1}} \int_0^{\min(r, c(v))}
\mathcal{A}(t, v)\, dt\, d\sigma(v)$,
where $\mathcal{A}(t, v)$ is the Jacobian of $\exp_p$ and $c(v)$ the cut distance.
By the Jacobi equation and the bound $\operatorname{Ric} \geq (n-1)\kappa$, one
shows that $\mathcal{A}(t,v)/\mathcal{A}_\kappa(t)$ is nonincreasing in $t$,
where $\mathcal{A}_\kappa(t) = \operatorname{sn}_\kappa(t)^{n-1}$ is the model
Jacobian.
\end{proof}

\begin{corollary}[Volume comparison]
Under the hypotheses of the theorem, for $0 < r \leq R$:
\[
    \frac{\operatorname{Vol}(B(p, R))}{\operatorname{Vol}(B(p, r))}
    \leq \frac{V_\kappa(R)}{V_\kappa(r)}.
\]
\end{corollary}

\begin{corollary}[Volume growth]
If $\operatorname{Ric} \geq 0$, then
$\operatorname{Vol}(B(p, r)) \leq \omega_n r^n$ (Euclidean volume).
\end{corollary}

\section{Application: Cheng's theorem}

\begin{theorem}[Cheng]
Let $(M^n, g)$ be complete with $\operatorname{Ric} \geq (n-1)\kappa > 0$.
If $\operatorname{diam}(M) = \pi/\sqrt{\kappa}$ (equality in Bonnet--Myers),
then $M$ is isometric to $S^n(1/\sqrt{\kappa})$.
\end{theorem}

\section{Laplacian comparison theorem}

\begin{theorem}[Laplacian comparison]
Let $(M^n, g)$ satisfy $\operatorname{Ric} \geq (n-1)\kappa$. The distance function
$r(x) = d(p, x)$ satisfies, in the sense of distributions:
\[
    \Delta r \leq (n-1)\frac{\operatorname{cn}_\kappa(r)}{\operatorname{sn}_\kappa(r)}.
\]
For $\kappa = 0$: $\Delta r \leq \frac{n-1}{r}$.
For $\kappa = 1$: $\Delta r \leq (n-1)\cot r$.
For $\kappa = -1$: $\Delta r \leq (n-1)\coth r$.
\end{theorem}

\section{Gromov's theorem on generators}

\begin{theorem}[Gromov]
If $(M^n, g)$ is complete with $\operatorname{Ric} \geq 0$, then the fundamental
group $\pi_1(M)$ is generated by at most $C(n)$ elements, where $C(n)$ depends
only on the dimension.
\end{theorem}

\section{Gromov--Hausdorff space}

\begin{definition}[Gromov--Hausdorff distance]
The \emph{Gromov--Hausdorff distance} between two compact metric spaces
$(X, d_X)$ and $(Y, d_Y)$ is:
\[
    d_{GH}(X, Y) = \inf \{ d_H(\varphi(X), \psi(Y)) \},
\]
where the infimum is over all isometric embeddings $\varphi : X \to Z$,
$\psi : Y \to Z$ into a metric space $(Z, d_Z)$, and $d_H$ is the Hausdorff
distance.
\end{definition}

\begin{theorem}[Gromov's precompactness]
The set of Riemannian manifolds $(M^n, g)$ satisfying
$\operatorname{Ric} \geq (n-1)\kappa$ and $\operatorname{diam}(M) \leq D$
is precompact in the Gromov--Hausdorff topology.
\end{theorem}

\section{Exercises}

\begin{exercise}
Verify the Bishop--Gromov theorem for balls in $S^n$ of curvature $1$ by
explicitly computing $\operatorname{Vol}(B(p, r))$ and comparing with $V_1(r)$.
\end{exercise}

\begin{exercise}
Show that if $(M, g)$ is complete with $K \geq 1$, then $\operatorname{diam}(M) \leq \pi$.
\end{exercise}

\begin{exercise}
Apply Toponogov's theorem to show that if $K \geq 0$ and $M$ is noncompact,
then $M$ contains a geodesic line (splitting theorem).
\end{exercise}

\begin{exercise}
Use volume comparison to show that if $\operatorname{Ric} \geq 0$ and
$\operatorname{Vol}(B(p,r)) \geq c\, r^n$ for some $c > 0$ and all $r$,
then $M$ has at most finitely many ends.
\end{exercise}

\begin{exercise}
Compute the Laplacian comparison for the distance function on $\mathbb{H}^3$
and verify that $\Delta r = 2\coth r$.
\end{exercise}

\begin{exercise}
Show that if $K \leq 0$ and $M$ is simply connected, then there are no conjugate
points (using the Rauch comparison theorem).
\end{exercise}
